\scanpage{116}
l'homomorphisme $\handwritten{\pi}_1^t(D-\overset{\handwritten{\scriptscriptstyle\vee}}{D})\handwritten{\rightarrow}\handwritten{\pi}_1^t(\unsure{\mathbb{P}|\deleted{R}P|\deleted{P}R}-\overset{\handwritten{\scriptscriptstyle\vee}}{X})$, et donne le théorème de conjugaison des cycles évanescents (donc le théorème d'irréductibilité de la cohomologie évanescente) pourvu qu'on ne soit pas dans le cas $\operatorname{car}.k=2$ et $\dim X$ impaire (NB \deleted{\unsure{ce|le} n'est pas,} dans la situation du th. de Griffiths, il n'y a donc pas d'ennui). Dans ce dernier cas, par contre, je ne sais \deleted{rien} pas prouver\inserted{la conjuga\unsure{\deleted{n}i|i}son sous $\handwritten{\pi}_1(D-\overset{\handwritten{\scriptscriptstyle\vee}}{D})$}, sauf pour un pinceau \uline{générique}, puisqu'alors la représentation envisagée est sauvage; \inserted{j'ignore si \unsure{ella|elle|\deleted{cela}\handwritten{elle}} reste vrai\unsure{\handwritten{e}|\deleted{e}} pour un pinceau \uline{général}.} Noter cependant que la conjugaison des sous-groupes de monodromie locale\deleted{\unsure{,|s}} \uline{dans} $\operatorname{Aut}(H^{n-1}\handwritten{(Y_{\bar{\eta}}))}$ (qui est très élémentaire) reste vraie, bien entendu, et celle-ci suffit pour prouver le théorème que j'ai proposé (au par. 6 de XVIII) caractérisant le cas où la condition (A) n'est \uline{pas} vérifiée.
% archive sheet number: 115
